\documentclass[12pt]{article}

\usepackage[spanish]{babel}
\usepackage[utf8]{inputenc}
\usepackage{amsmath}
\usepackage{geometry}
\usepackage{booktabs}

\geometry{margin=1in}

\title{Módulo 2: Ejercicios Prácticos en R\\
Identificación y Medición de Riesgos Previsionales}
\author{MACROPOL}
\date{}

\begin{document}

\maketitle

\section{Contexto}

El objetivo de estos ejercicios es traducir riesgos financieros en métricas previsionales relevantes. Se parte de la medición de riesgo de mercado y se avanza hacia la evaluación del impacto sobre la tasa de reemplazo y la suficiencia de pensiones.

\section{Ejercicio 1: Value at Risk (VaR) y Expected Shortfall (ES)}

\subsection{Objetivo}
Medir el riesgo de mercado de un portafolio.

\subsection{Código en R}
\begin{verbatim}
set.seed(123)

returns <- rnorm(1000, mean = 0.0005, sd = 0.02)

VaR_95 <- quantile(returns, 0.05)
ES_95 <- mean(returns[returns < VaR_95])

VaR_95
ES_95
\end{verbatim}

\subsection{Preguntas}
\begin{itemize}
\item ¿Qué representa el VaR?
\item ¿Por qué el ES es más informativo en escenarios extremos?
\end{itemize}

\subsection{Interpretación}
El VaR mide un umbral de pérdida, mientras que el ES captura la severidad de pérdidas en crisis.

\section{Ejercicio 2: Dinámica del Fondo de Pensiones}

\subsection{Modelo}
\begin{equation}
W_{t+1} = W_t (1 + r_t) + c_t
\end{equation}

\subsection{Código en R}
\begin{verbatim}
T <- 30
W <- numeric(T)
W[1] <- 10000
c <- 1000

returns <- rnorm(T, 0.05, 0.15)

for(t in 2:T){
  W[t] <- W[t-1]*(1 + returns[t]) + c
}

plot(W, type="l")
\end{verbatim}

\subsection{Preguntas}
\begin{itemize}
\item ¿Cómo afecta la volatilidad la trayectoria del fondo?
\item Compare con un escenario de retornos negativos.
\end{itemize}

\section{Ejercicio 3: Tasa de Reemplazo}

\subsection{Modelo}
\begin{equation}
RR = \frac{W_T}{\text{factor actuarial} \times salario}
\end{equation}

\subsection{Código en R}
\begin{verbatim}
final_wealth <- W[T]
salary <- 20000
annuity_factor <- 15

RR <- final_wealth / (annuity_factor * salary)
RR
\end{verbatim}

\subsection{Pregunta}
¿Es la tasa de reemplazo suficiente (por ejemplo, mayor a 40\%)?

\section{Ejercicio 4: Probabilidad de Insuficiencia}

\subsection{Objetivo}
Calcular la probabilidad de no alcanzar una pensión mínima.

\subsection{Código en R}
\begin{verbatim}
set.seed(123)

N <- 5000
RR_sim <- numeric(N)

for(i in 1:N){
  W <- 10000
  
  for(t in 1:30){
    r <- rnorm(1, 0.05, 0.15)
    W <- W*(1 + r) + 1000
  }
  
  RR_sim[i] <- W / (15 * 20000)
}

prob_insuf <- mean(RR_sim < 0.4)
prob_insuf
\end{verbatim}

\subsection{Pregunta}
¿Qué proporción de afiliados no alcanza una pensión adecuada?

\section{Ejercicio 5: Stress Test Integrado}

\subsection{Escenario Adverso}
\begin{itemize}
\item Retornos negativos
\item Mayor volatilidad
\item Aumento de longevidad
\end{itemize}

\subsection{Código en R}
\begin{verbatim}
set.seed(123)

N <- 5000
RR_stress <- numeric(N)

for(i in 1:N){
  W <- 10000
  
  for(t in 1:30){
    r <- rnorm(1, -0.02, 0.20)
    W <- W*(1 + r) + 1000
  }
  
  RR_stress[i] <- W / (20 * 20000)
}

mean(RR_stress)
mean(RR_stress < 0.4)
\end{verbatim}

\section{Análisis Comparativo}

\begin{center}
\begin{tabular}{lcc}
\toprule
Métrica & Baseline & Adverso \\
\midrule
Tasa de reemplazo &  &  \\
Probabilidad de insuficiencia &  &  \\
\bottomrule
\end{tabular}
\end{center}

\section{Interpretación}

\begin{itemize}
\item El VaR mide riesgo financiero, pero no bienestar previsional.
\item El ES captura pérdidas en crisis.
\item La tasa de reemplazo conecta con el resultado final del sistema.
\item La probabilidad de insuficiencia mide el impacto social del riesgo.
\end{itemize}

\section{Mensaje Final}

\begin{center}
\textit{El riesgo no importa por sí mismo, sino por cómo afecta la pensión final.}
\end{center}

\end{document}